% Compile with LuaLaTeX (recommended) or XeLaTeX
\documentclass[11pt,a4paper]{article}

% ==========================================
% 1. MODERN FONTS & LANGUAGES
% ==========================================
\usepackage{fontspec}
% Tip: For Romance languages and IPA, "Charis SIL" or "Doulos SIL" are the best fonts. 
% If you have them installed, uncomment the next line:
% \setmainfont{Charis SIL} 

\usepackage{polyglossia}
\setmainlanguage{english} % change to german if needed
% Load the Romance languages you study
\setotherlanguages{french,spanish,italian,portuguese,romanian,latin}

% ==========================================
% 1b. TYPOGRAPHY & QUOTES
% ==========================================
\usepackage{microtype}   % Better line breaks, especially helpful for multilingual docs
\usepackage{csquotes}    % Language-aware quotes: \enquote{...} picks the right marks
\usepackage{booktabs}
\usepackage{graphicx}

% ==========================================
% 2. LINGUISTICS & PHONETICS
% ==========================================
\usepackage{tipauni}     % Modern, LuaLaTeX-safe IPA symbols
\usepackage{gb4e}        % For interlinear glosses (standard in linguistics)
\noautomath              % CRITICAL: Prevents gb4e from crashing with modern packages
\usepackage{forest}      % For drawing modern syntax trees

% ==========================================
% 3. MATH & LAYOUT
% ==========================================
\usepackage{amsmath,amssymb,mathtools}
\usepackage{geometry}
\geometry{margin=2.5cm}
\usepackage{hyperref}
\usepackage{xcolor}
\usepackage{siunitx}

% ---------- Bibliography (uncomment when you have a .bib file) ----------
% \usepackage[style=authoryear,backend=biber]{biblatex}
% \addbibresource{references.bib}

\title{Computational Linguistics Starter Document}
\author{Your Name}
\date{\today}

\begin{document}

\maketitle

\section{Simple Statistical Formula}

Conditional probability:
\begin{equation}
  P(w_i \mid w_{i-1}) = \frac{C(w_{i-1}, w_i)}{C(w_{i-1})}
\end{equation}

Precision and Recall:
\begin{align*}
  \text{Precision} &= \frac{TP}{TP + FP} \\
  \text{Recall}    &= \frac{TP}{TP + FN}
\end{align*}

\section{Linguistic Example with Gloss}

With \texttt{csquotes}, write \verb|\enquote{...}| and let LaTeX pick the
right quotation marks for the current language: \enquote{So she said.}

\begin{exe}
  \ex \gll  Elle a mangé une pomme. \\
            She has eaten an  apple. \\
      \glt  `She ate an apple.'
\end{exe}

\section{IPA Example}

English ``time'': \textipa{[aI]} \\
French \textfrench{pain}: \textipa{[pE\~{}]}

\section{Simple Syntax Tree}

\begin{center}
\begin{forest}
  [S
    [NP [Det the] [N cat]]
    [VP [V sleeps]]
  ]
\end{forest}
\end{center}

\section{Multilingual Text}

English text.
\textfrench{Voici un exemple en français.}
\textspanish{Aquí un ejemplo en español.}
\textitalian{Ecco un esempio in italiano.}

\section{Technical Notes}

If you prefer to avoid \texttt{\textbackslash noautomath} entirely, there are two modern paths:
\begin{itemize}
    \item Drop \texttt{tipa}/\texttt{tipauni} and type IPA characters directly as Unicode, with a font that covers them (Charis SIL, Doulos SIL, Noto Sans).
    \item Switch from \texttt{gb4e} to \texttt{expex} or \texttt{linguex}, both of which have fewer conflicts with other packages.
\end{itemize}
For now, \texttt{tipauni} + \texttt{gb4e} + \texttt{\textbackslash noautomath} is a perfectly fine and widely used combination.

% \printbibliography

\end{document}
